\subsection{COMMUTATORS}

DEFINITION 3. Let G be a group and x and y two elements of G. The element \(x^{-1}y^{-1}xy\) of G is called the commutator of x and y.

\((x,y)\) is used to denote the commutator of x and y. Then obviously

\[
(y, x) = (x, y) ^ {- 1}.
\]

For \(x\) and \(y\) to commute it is necessary and sufficient that \((x, y) = e\) . More generally,

\[
x y = y x (x, y).
\]

On the other hand we write

\[
x ^ {y} = y ^ {- 1} x y = x (x, y) = (y, x ^ {- 1}) x.\tag{3}
\]

As the mapping \(x \mapsto x^y\) is the inner automorphism \(\operatorname{Int}(y^{-1})\) , \((x, y)^z = (x^z, y^z)\) for all \(x, y, z \in G\) .

For \(x, y, z \in G\) , we prove the following relations:

\[
(x, y z) = (x, z). (x, y) ^ {z} = (x, z). (z, (y, x)). (x, y)\tag{4}
\]

(4 bis) \((xy,z) = (x,z)^y.(y,z) = (x,z).(((x,z),y).(y,z))\)

\begin{enumerate}
\def\labelenumi{(\arabic{enumi})}
\setcounter{enumi}{4}
\item
  \((x^{y}, (y, z)).(y^{z}, (z, x)).(z^{x}, (x, y)) = e\)
\item
  \((x, yz).(z, xy).(y, zx) = e\)
\end{enumerate}

(6 bis) \((xy, z). (yz, x). (zx, y) = e.\)

Now

\[
\begin{array}{r l} (x, y z) & = x ^ {- 1} z ^ {- 1} y ^ {- 1} x y z = (x, z) z ^ {- 1} x ^ {- 1} y ^ {- 1} x y z = (x, z) (x, y) ^ {z} \\ & = (x, z) (z, (x, y) ^ {- 1}) (x, y) \end{array}
\]

by (3), which proves (4). Formula (4 bis) follows similarly. On the other hand,

\[
\begin{array}{r l} (x ^ {y}, (y, z)) & = (x ^ {y}) ^ {- 1} (z, y) (x ^ {y}) (y, z) \\ & = y ^ {- 1} x ^ {- 1} y z ^ {- 1} y ^ {- 1} z y y ^ {- 1} x y y ^ {- 1} z ^ {- 1} y z \\ & = (y z y ^ {- 1} x y) ^ {- 1} (z x z ^ {- 1} y z). \end{array}
\]

Then writing \(u = yzy^{-1}xy\) , \(v = zxz^{-1}yz\) and \(w = xyx^{-1}zx\) , we obtain

\[
(x ^ {y}, (y, z)) = u ^ {- 1} v.
\]

By cyclically permuting \(x, y, z\) , we deduce \((y^z, (z, x)) = v^{-1}w\) and

\[
\left(z ^ {x}, (x, y)\right) = w ^ {- 1} u,
\]

which immediately imply (5). Finally, (6) follows by multiplying together the two sides in the three formulae obtained by cyclically permuting x, y, z in the formula \((x, yz) = x^{-1}z^{-1}y^{-1}xyz = (yzx)^{-1}(xyz)\) , and similarly for (6 bis).

If A and B are two subgroups of G, (A, B) denotes the subgroup generated by the commutators \((a, b)\) with \(a \in A\) and \(b \in B\) . Then \((A, B) = \{e\}\) if and only if A centralizes B. \((A, B) \subset A\) if and only if B normalizes A. If A and B are normal (resp. characteristic), so is \((A, B)\) .

PROPOSITION 5. Let A, B, C be three subgroups of G.

\begin{enumerate}
\def\labelenumi{(\roman{enumi})}
\item
  The subgroup A normalizes the subgroup (A, B).
\item
  If the subgroup (B, C) normalizes A, the subgroup (A, (B, C)) is generated by the elements \((a, (b, c))\) with \(a \in \mathbf{A}\) , \(b \in \mathbf{B}\) and \(c \in \mathbf{C}\) .
\item
  If A, B and C are normal, then
\end{enumerate}

\[
(\mathrm{A}, (\mathrm{B}, \mathrm{C})) \subset (\mathrm{C}, (\mathrm{B}, \mathrm{A})). (\mathrm{B}, (\mathrm{C}, \mathrm{A})).
\]

By (4 bis), for \(a, a' \in \mathbf{A}\) and \(b \in \mathbf{B}\) ,

\[
(a, b) ^ {a ^ {\prime}} = (a a ^ {\prime}, b). (a ^ {\prime}, b) ^ {- 1},
\]

whence (i). Suppose now that (B, C) normalizes A. For \(a \in \mathbf{A}\) , \(b \in \mathbf{B}\) , \(c \in \mathbf{C}\) and \(x \in \mathbf{G}\) , (4) implies

\[
(a, (b, c). x) = (a, x). (x, ((b, c), a)) (a, (b, c))
\]

and \(((b, c), a) \in \mathbf{A}\) since (B, C) normalizes A, whence by induction on \(p\) the fact that \(\left(a, \prod_{i=1}^{p} (b_i, c_i)\right)\) , for \(b_i \in \mathbf{B}\) , \(c_i \in \mathbf{C}\) , belongs to the subgroup generated by the elements of the form \((a, (b, c))\) . If finally A, B and C are normal, so are the subgroups (A, (B, C)), (C, (B, A)) and (B, (C, A)). It therefore suffices by (ii) to show that

\[
(a, (b, c)) \in (\mathrm{C}, (\mathrm{B}, \mathrm{A})). (\mathrm{B}, (\mathrm{C}, \mathrm{A}))
\]

for all \(a \in \mathbf{A}\) , \(b \in \mathbf{B}\) and \(c \in \mathbf{C}\) . Now by (5), writing \(a^{b^{-1}} = u\)

\[
(a, (b, c)) = ((u ^ {b}), (b, c)) = (c ^ {u}, (u, b)) ^ {- 1}. (b ^ {c}, (c, u)) ^ {- 1},
\]

whence (iii).

DEFINITION 4. Let G be a group. The subgroup generated by the commutators of elements of G is called the derived group of G.

The derived group of G is thus the subgroup (G, G). It is also denoted by D(G). By an abuse of language, it is sometimes called the commutator group of G although it is in general distinct from the set of commutators of elements of G (Exercise 16). D(G) = \{e\} if and only if G is commutative.

PROPOSITION 6. Let \(f: \mathbf{G} \to \mathbf{G}'\) be a group homomorphism. Then \(f(\mathbf{D}(\mathbf{G})) \subset \mathbf{D}(\mathbf{G}')\) . If \(f\) is surjective, the homomorphism of \(\mathbf{D}(\mathbf{G})\) into \(\mathbf{D}(\mathbf{G}')\) the restriction of \(f\) is surjective.

The image under \(f\) of a commutator of elements of \(G\) is a commutator of elements of \(G'\) . If \(f\) is surjective, the image under \(f\) of the set of commutators of \(G\) is the set of commutators of \(G'\) . The proposition thus follows from § 4, no. 3, Corollary 3 to Proposition 2.

COROLLARY 1. The derived group of a group G is a characteristic subgroup of G. In particular it is a normal subgroup of G.

COROLLARY 2. Let G be a group. The quotient group G/D(G) is commutative. Let \(\pi: G \to G/D(G)\) be the canonical homomorphism. Every homomorphism \(f\) of G into a commutative group \(G'\) can be expressed uniquely in the form \(f = \bar{f} \circ \pi\) , where \(\bar{f}: G/D(G) \to G'\) is a homomorphism.

Now \(\pi (\mathbf{D}(\mathbf{G})) = \{e\}\) . As \(\pi\) is surjective, it follows that \(\mathbf{D}(\mathbf{G} / \mathbf{D}(\mathbf{G})) = \{e\}\) , whence the first assertion. The second follows from § 4, no. 4, Proposition 5.

COROLLARY 3. Let H be a subgroup of G. The following conditions are equivalent:

\begin{enumerate}
\def\labelenumi{(\roman{enumi})}
\item
  \(\mathbf{H} \supset \mathbf{D}(\mathbf{G})\) ;
\item
  H is a normal subgroup and G/H is commutative.
\item
  (ii) \(\Rightarrow\) (i) by Corollary 2 and (i) \(\Rightarrow\) (ii) by §4, no. 7, Theorem 4, since every subgroup of a commutative group is normal.
\end{enumerate}

COROLLARY 4. Let G be a group and X a subset of G which generates G. The group D(G) is the normal subgroup of G generated by the commutators of elements of X.

Let H be the normal subgroup of G generated by the commutators of elements of X and \(\phi: G \to G/H\) the canonical homomorphism. The set \(\phi(X)\) generates G/H. The elements of \(\phi(X)\) are pairwise permutable and hence H is commutative ( \(\S 4\) , no. 3, Corollary 2 to Proposition 2). Hence (Corollary 3) H contains D(G). On the other hand, obviously \(H \subset D(G)\) .

Remarks. (1) Corollary 2 can also be expressed by saying that G/D(G), together with \(\pi\) , is a solution of the universal mapping problem for G, relative to commutative groups and homomorphisms from G to commutative groups.

\begin{enumerate}
\def\labelenumi{(\arabic{enumi})}
\setcounter{enumi}{1}
\tightlist
\item
  Under the hypotheses of Corollary 4, the subgroup generated by the commutators of elements of X is contained in D(G) but is not in general equal to D(G) (cf.~Exercise 15e).
\end{enumerate}

Examples. (1) If \(\mathbf{G}\) is a non-commutative simple group, then \(\mathrm{D}(\mathbf{G}) = \mathbf{G}\) . Therefore every homomorphism of \(\mathbf{G}\) into a commutative group is trivial.

\begin{enumerate}
\def\labelenumi{(\arabic{enumi})}
\setcounter{enumi}{1}
\tightlist
\item
  The derived group of the symmetric group \(\mathfrak{S}_n\) is the alternating group \(\mathfrak{A}_n\) . For \(\mathfrak{A}_n\) is generated by the products of two transpositions; if \(\tau = \tau_{x,y}\) and \(\tau' = \tau_{x',y'}\) are two transpositions, let \(\sigma\) be a permutation such that \(\sigma(x') = x\) and \(\sigma(y') = y\) . Then \(\tau' = \sigma^{-1}\tau\sigma\) and \(\tau\tau' = \tau^{-1}\tau' = \tau^{-1}\sigma^{-1}\tau\sigma\) is a commutator. Hence \(\mathfrak{A}_n \subset D(\mathfrak{S}_n)\) . As \(\mathfrak{S}_n / \mathfrak{A}_n\) is commutative, \(\mathfrak{A}_n \supset D(\mathfrak{S}_n)\) (Corollary 3).
\end{enumerate}

